\documentclass[10pt,a4paper]{amsart}
\usepackage{amsmath,amssymb,amsthm,graphicx}
% Retained minimal preamble for the Stacks Project chapter more-morphisms.tex (More on Morphisms).
% Verbatim excerpt of the chapter preamble: Xy-pic setup (native preamble lines 14-18) and the native
% theorem declarations with the mathematical macros (native preamble lines 154-217). Copied mechanically;
% nothing added, removed or redefined. The wrapper class and the amsmath/amssymb/amsthm/graphicx base
% packages are supplied by the assembler, so the class/externaldocument part of the native preamble is omitted.
% For commutative diagrams we use Xy-pic
\usepackage[all]{xy}

% We use 2cell for 2-commutative diagrams.
\xyoption{2cell}
% Theorem environments.
%
\theoremstyle{plain}
\newtheorem{theorem}[subsection]{Theorem}
\newtheorem{proposition}[subsection]{Proposition}
\newtheorem{lemma}[subsection]{Lemma}

\theoremstyle{definition}
\newtheorem{definition}[subsection]{Definition}
\newtheorem{example}[subsection]{Example}
\newtheorem{exercise}[subsection]{Exercise}
\newtheorem{situation}[subsection]{Situation}

\theoremstyle{remark}
\newtheorem{remark}[subsection]{Remark}
\newtheorem{remarks}[subsection]{Remarks}

\numberwithin{equation}{subsection}

% Macros
%
\def\lim{\mathop{\mathrm{lim}}\nolimits}
\def\colim{\mathop{\mathrm{colim}}\nolimits}
\def\Spec{\mathop{\mathrm{Spec}}}
\def\Hom{\mathop{\mathrm{Hom}}\nolimits}
\def\Ext{\mathop{\mathrm{Ext}}\nolimits}
\def\SheafHom{\mathop{\mathcal{H}\!\mathit{om}}\nolimits}
\def\SheafExt{\mathop{\mathcal{E}\!\mathit{xt}}\nolimits}
\def\Sch{\mathit{Sch}}
\def\Mor{\mathop{\mathrm{Mor}}\nolimits}
\def\Ob{\mathop{\mathrm{Ob}}\nolimits}
\def\Sh{\mathop{\mathit{Sh}}\nolimits}
\def\NL{\mathop{N\!L}\nolimits}
\def\CH{\mathop{\mathrm{CH}}\nolimits}
\def\proetale{{pro\text{-}\acute{e}tale}}
\def\etale{{\acute{e}tale}}
\def\QCoh{\mathit{QCoh}}
\def\Ker{\mathop{\mathrm{Ker}}}
\def\Im{\mathop{\mathrm{Im}}}
\def\Coker{\mathop{\mathrm{Coker}}}
\def\Coim{\mathop{\mathrm{Coim}}}

% Boxtimes
%
\DeclareMathSymbol{\boxtimes}{\mathbin}{AMSa}{"02}

%
% Macros for moduli stacks/spaces
%
\def\QCohstack{\mathcal{QC}\!\mathit{oh}}
\def\Cohstack{\mathcal{C}\!\mathit{oh}}
\def\Spacesstack{\mathcal{S}\!\mathit{paces}}
\def\Quotfunctor{\mathrm{Quot}}
\def\Hilbfunctor{\mathrm{Hilb}}
\def\Curvesstack{\mathcal{C}\!\mathit{urves}}
\def\Polarizedstack{\mathcal{P}\!\mathit{olarized}}
\def\Complexesstack{\mathcal{C}\!\mathit{omplexes}}
% \Pic is the operator that assigns to X its picard group, usage \Pic(X)
% \Picardstack_{X/B} denotes the Picard stack of X over B
% \Picardfunctor_{X/B} denotes the Picard functor of X over B
\def\Pic{\mathop{\mathrm{Pic}}\nolimits}
\def\Picardstack{\mathcal{P}\!\mathit{ic}}
\def\Picardfunctor{\mathrm{Pic}}
\def\Deformationcategory{\mathcal{D}\!\mathit{ef}}
\newenvironment{reference}{}{}
\newenvironment{slogan}{}{}
\newenvironment{history}{}{}
\usepackage{cleveref}
\crefname{originallemma}{original lemma}{original lemmas}
\Crefname{originallemma}{Original lemma}{Original lemmas}
\makeatletter
\newlabel{constructions-lemma-projective-space-bundle}{{01OE}{source}{}{sourcepointer-1}{}}
\newlabel{constructions-lemma-projective-space-bundle@cref}{{[originallemma][0][]01OE}{[1][0][]source}{}{}{}}
\makeatother
\crefname{originallemma}{original lemma}{original lemmas}
\Crefname{originallemma}{Original lemma}{Original lemmas}
\makeatletter
\newlabel{divisors-lemma-closed-subscheme-proj}{{0801}{source}{}{sourcepointer-2}{}}
\newlabel{divisors-lemma-closed-subscheme-proj@cref}{{[originallemma][0][]0801}{[1][0][]source}{}{}{}}
\makeatother
\crefname{originallemma}{original lemma}{original lemmas}
\Crefname{originallemma}{Original lemma}{Original lemmas}
\makeatletter
\newlabel{divisors-lemma-relative-proj-projective}{{0B3U}{source}{}{sourcepointer-3}{}}
\newlabel{divisors-lemma-relative-proj-projective@cref}{{[originallemma][0][]0B3U}{[1][0][]source}{}{}{}}
\makeatother
\crefname{originallemma}{original lemma}{original lemmas}
\Crefname{originallemma}{Original lemma}{Original lemmas}
\makeatletter
\newlabel{lemma-quasi-projective}{{0B42}{source}{}{sourcepointer-4}{}}
\newlabel{lemma-quasi-projective@cref}{{[originallemma][0][]0B42}{[1][0][]source}{}{}{}}
\makeatother
\crefname{originallemma}{original lemma}{original lemmas}
\Crefname{originallemma}{Original lemma}{Original lemmas}
\makeatletter
\newlabel{morphisms-definition-projective}{{01W8}{source}{}{sourcepointer-5}{}}
\newlabel{morphisms-definition-projective@cref}{{[originallemma][0][]01W8}{[1][0][]source}{}{}{}}
\makeatother
\crefname{originallemma}{original lemma}{original lemmas}
\Crefname{originallemma}{Original lemma}{Original lemmas}
\makeatletter
\newlabel{morphisms-definition-quasi-projective}{{01VW}{source}{}{sourcepointer-6}{}}
\newlabel{morphisms-definition-quasi-projective@cref}{{[originallemma][0][]01VW}{[1][0][]source}{}{}{}}
\makeatother
\crefname{originallemma}{original lemma}{original lemmas}
\Crefname{originallemma}{Original lemma}{Original lemmas}
\makeatletter
\newlabel{morphisms-definition-very-ample}{{01VM}{source}{}{sourcepointer-7}{}}
\newlabel{morphisms-definition-very-ample@cref}{{[originallemma][0][]01VM}{[1][0][]source}{}{}{}}
\makeatother
\crefname{originallemma}{original lemma}{original lemmas}
\Crefname{originallemma}{Original lemma}{Original lemmas}
\makeatletter
\newlabel{morphisms-lemma-H-projective}{{01W9}{source}{}{sourcepointer-8}{}}
\newlabel{morphisms-lemma-H-projective@cref}{{[originallemma][0][]01W9}{[1][0][]source}{}{}{}}
\makeatother
\crefname{originallemma}{original lemma}{original lemmas}
\Crefname{originallemma}{Original lemma}{Original lemmas}
\makeatletter
\newlabel{morphisms-lemma-locally-projective-proper}{{01WC}{source}{}{sourcepointer-9}{}}
\newlabel{morphisms-lemma-locally-projective-proper@cref}{{[originallemma][0][]01WC}{[1][0][]source}{}{}{}}
\makeatother
\crefname{originallemma}{original lemma}{original lemmas}
\Crefname{originallemma}{Original lemma}{Original lemmas}
\makeatletter
\newlabel{morphisms-lemma-projective-is-quasi-projective-proper}{{0BCL}{source}{}{sourcepointer-10}{}}
\newlabel{morphisms-lemma-projective-is-quasi-projective-proper@cref}{{[originallemma][0][]0BCL}{[1][0][]source}{}{}{}}
\makeatother
\crefname{originallemma}{original lemma}{original lemmas}
\Crefname{originallemma}{Original lemma}{Original lemmas}
\makeatletter
\newlabel{morphisms-lemma-projective-over-quasi-projective-is-H-projective}{{087S}{source}{}{sourcepointer-11}{}}
\newlabel{morphisms-lemma-projective-over-quasi-projective-is-H-projective@cref}{{[originallemma][0][]087S}{[1][0][]source}{}{}{}}
\makeatother
\title{Eight equivalent descriptions of projective morphisms over a base with an ample invertible sheaf}
\author{The Stacks Project Authors}
\date{stacks-project:more-morphisms.tex:local-snapshot-20261011; GNU FDL 1.2 or later}
\begin{document}
\maketitle
\noindent\textit{Editorial note.} Statement span is lines 14027-14045 of sources/stacks-more-morphisms/pinned-original-source.tex; proof span is lines 14047-14100, the complete author proof (52 body lines).\par\medskip
\noindent\textit{Editorial note.} Navigation stubs: 11 cited-but-not-included labels are supplied as non-load-bearing pointers with their pinned tags where the snapshot tag index has them.\par\medskip
\noindent\textit{Editorial note.} extra\_packages is empty on purpose: the retained preamble already loads xy with options and a second declaration clashes in the assembler wrapper.\par\medskip
\noindent\textit{Editorial note.} Text is retained verbatim, including the author's own typos and placeholders.\par\medskip
\noindent\textit{Editorial note.} Hub lemma of the chapter's projective/quasi-projective series; complements L16072 (quasi-projective) without overlap.\par\medskip
\noindent\textit{Allowed standard prerequisites (editorial).}\begin{itemize}
\item Explicit prerequisite boundary. The quasi-projective characterization card (same chapter, collected as L16072), the Morphisms chapter dictionary between projective/H-projective/quasi-projective notions, and the relative Proj machinery (Divisors, lemmas relative-proj-projective and closed-subscheme-proj; Constructions, lemma-projective-space-bundle).
\end{itemize}\medskip
\ifdefined\hypertarget\hypertarget{sourcepointer-1}{}\fi
\noindent\textit{Source reference.} Original lemma 01OE: sources/stacks-more-morphisms/pinned-original-source.tex, source lines 14100--14100. Cited by the retained statement or proof but not part of the collected spans; supplied only as a navigation pointer. The argument does not depend on its content.\par\medskip
\ifdefined\hypertarget\hypertarget{sourcepointer-2}{}\fi
\noindent\textit{Source reference.} Original lemma 0801: sources/stacks-more-morphisms/pinned-original-source.tex, source lines 14100--14100. Cited by the retained statement or proof but not part of the collected spans; supplied only as a navigation pointer. The argument does not depend on its content.\par\medskip
\ifdefined\hypertarget\hypertarget{sourcepointer-3}{}\fi
\noindent\textit{Source reference.} Original lemma 0B3U: sources/stacks-more-morphisms/pinned-original-source.tex, source lines 14100--14100. Cited by the retained statement or proof but not part of the collected spans; supplied only as a navigation pointer. The argument does not depend on its content.\par\medskip
\ifdefined\hypertarget\hypertarget{sourcepointer-4}{}\fi
\noindent\textit{Source reference.} Original lemma 0B42: sources/stacks-more-morphisms/pinned-original-source.tex, source lines 13840--13853. Cited by the retained statement or proof but not part of the collected spans; supplied only as a navigation pointer. The argument does not depend on its content.\par\medskip
\ifdefined\hypertarget\hypertarget{sourcepointer-5}{}\fi
\noindent\textit{Source reference.} Original lemma 01W8: sources/stacks-more-morphisms/pinned-original-source.tex, source lines 14100--14100. Cited by the retained statement or proof but not part of the collected spans; supplied only as a navigation pointer. The argument does not depend on its content.\par\medskip
\ifdefined\hypertarget\hypertarget{sourcepointer-6}{}\fi
\noindent\textit{Source reference.} Original lemma 01VW: sources/stacks-more-morphisms/pinned-original-source.tex, source lines 14100--14100. Cited by the retained statement or proof but not part of the collected spans; supplied only as a navigation pointer. The argument does not depend on its content.\par\medskip
\ifdefined\hypertarget\hypertarget{sourcepointer-7}{}\fi
\noindent\textit{Source reference.} Original lemma 01VM: sources/stacks-more-morphisms/pinned-original-source.tex, source lines 14100--14100. Cited by the retained statement or proof but not part of the collected spans; supplied only as a navigation pointer. The argument does not depend on its content.\par\medskip
\ifdefined\hypertarget\hypertarget{sourcepointer-8}{}\fi
\noindent\textit{Source reference.} Original lemma 01W9: sources/stacks-more-morphisms/pinned-original-source.tex, source lines 14100--14100. Cited by the retained statement or proof but not part of the collected spans; supplied only as a navigation pointer. The argument does not depend on its content.\par\medskip
\ifdefined\hypertarget\hypertarget{sourcepointer-9}{}\fi
\noindent\textit{Source reference.} Original lemma 01WC: sources/stacks-more-morphisms/pinned-original-source.tex, source lines 14100--14100. Cited by the retained statement or proof but not part of the collected spans; supplied only as a navigation pointer. The argument does not depend on its content.\par\medskip
\ifdefined\hypertarget\hypertarget{sourcepointer-10}{}\fi
\noindent\textit{Source reference.} Original lemma 0BCL: sources/stacks-more-morphisms/pinned-original-source.tex, source lines 14100--14100. Cited by the retained statement or proof but not part of the collected spans; supplied only as a navigation pointer. The argument does not depend on its content.\par\medskip
\ifdefined\hypertarget\hypertarget{sourcepointer-11}{}\fi
\noindent\textit{Source reference.} Original lemma 087S: sources/stacks-more-morphisms/pinned-original-source.tex, source lines 14100--14100. Cited by the retained statement or proof but not part of the collected spans; supplied only as a navigation pointer. The argument does not depend on its content.\par\medskip
\begin{lemma}
\label{lemma-projective}
Let $S$ be a scheme which has an ample invertible sheaf.
Let $f : X \to S$ be a morphism of schemes. The following are
equivalent
\begin{enumerate}
\item $X \to S$ is projective,
\item $X \to S$ is H-projective,
\item $X \to S$ is quasi-projective and proper,
\item $X \to S$ is H-quasi-projective and proper,
\item $X \to S$ is proper and $X$ has an ample invertible sheaf,
\item $X \to S$ is proper and there exists an $f$-ample invertible sheaf,
\item $X \to S$ is proper and there exists an $f$-very ample invertible sheaf,
\item there is a quasi-coherent graded $\mathcal{O}_S$-algebra $\mathcal{A}$
generated by $\mathcal{A}_1$ over $\mathcal{A}_0$ with $\mathcal{A}_1$ a
finite type $\mathcal{O}_S$-module such that
$X = \underline{\text{Proj}}_S(\mathcal{A})$. 
\end{enumerate}
\end{lemma}

\begin{proof}
Observe first that in each case the morphism $f$ is proper, see
Morphisms, Lemmas \ref{morphisms-lemma-H-projective} and
\ref{morphisms-lemma-locally-projective-proper}.
Hence it suffices to prove the equivalence of the notions in
case $f$ is a proper morphism. We will use this without further
mention in the following.

\medskip\noindent
The equivalences (1) $\Leftrightarrow$ (3) and
(2) $\Leftrightarrow$ (4) are
Morphisms, Lemma \ref{morphisms-lemma-projective-is-quasi-projective-proper}.

\medskip\noindent
The implication (2) $\Rightarrow$ (1) is
Morphisms, Lemma \ref{morphisms-lemma-H-projective}.

\medskip\noindent
The implications (1) $\Rightarrow$ (2) and (3) $\Rightarrow$ (4) are
Morphisms, Lemma
\ref{morphisms-lemma-projective-over-quasi-projective-is-H-projective}.

\medskip\noindent
The implication (1) $\Rightarrow$ (7) is immediate from
Morphisms, Definitions \ref{morphisms-definition-projective} and
\ref{morphisms-definition-very-ample}.

\medskip\noindent
The conditions (3) and (6) are equivalent by
Morphisms, Definition \ref{morphisms-definition-quasi-projective}.

\medskip\noindent
Thus (1) -- (4), (6) are equivalent and imply (7). By
Lemma \ref{lemma-quasi-projective}
conditions (3), (5), and (7) are equivalent.
Thus we see that (1) -- (7) are equivalent.

\medskip\noindent
By Divisors, Lemma \ref{divisors-lemma-relative-proj-projective}
we see that (8) implies (1). Conversely, if (2) holds, then
we can choose a closed immersion
$$
i :
X
\longrightarrow
\mathbf{P}^n_S = \underline{\text{Proj}}_S(\mathcal{O}_S[T_0, \ldots, T_n]).
$$
See Constructions, Lemma \ref{constructions-lemma-projective-space-bundle}
for the equality. By
Divisors, Lemma \ref{divisors-lemma-closed-subscheme-proj}
we see that $X$ is the relative Proj of a quasi-coherent graded quotient
algebra $\mathcal{A}$ of $\mathcal{O}_S[T_0, \ldots, T_n]$.
Then $\mathcal{A}$ satisfies the conditions of (8).
\end{proof}

\end{document}
